% GENERATED FILE. Edit canonical lessons, then run npm run generate:books.
% canonical-source-hash: fnv1a64:c10db42ef77a434d
% canonical-lessons: KA-C449-cakra, KA-C449-tayar, KA-C449-stepni, KA-C449-pankcar-agu, KA-C449-brek

\chapter{A Flat Tyre}
\label{ch:ka-a-flat-tyre}
\hlchaptermodality{449}

\begin{chapteropening}
\textbf{By the end of this chapter:} \emph{I can say a wheel, a tyre, a spare tyre, to get a puncture and a brake.}

Complete the last lesson of chapter 449: I can say a wheel, a tyre, a spare tyre, to get a puncture and a brake.
\end{chapteropening}

\section[\texorpdfstring{\kn{ಚಕ್ರ} (cakra)}{ಚಕ್ರ (cakra)}]{\kn{ಚಕ್ರ} (cakra) — a wheel}

\label{lesson:KA-C449-cakra}



\begin{quote}
Before the new one: say the Kannada for a battery, then the Kannada for to ride.
\end{quote}

\subsection*{You'll want to know: \kn{ಚಕ್ರ}}
\textbf{\kn{ಚಕ್ರ}} — \emph{cakra} — \textquotedblleft{}a wheel\textquotedblright{}.

\begin{quote}
\kn{ಸೈಕಲ್ನ} \kn{ಹಿಂದಿನ} \kn{ಚಕ್ರ} \kn{ಡೊಂಕಾಗಿದೆ}.
\end{quote}

\begin{quote}
\emph{saikalna hindina cakra ḍoṅkāgide.}
\end{quote}

\begin{quote}
The back wheel of the bicycle is bent.
\end{quote}

Plural: \kn{ಚಕ್ರಗಳು}, \emph{cakragaḷu}.

\kn{ದ್ವಿಚಕ್ರ} \kn{ವಾಹನ}, a two-wheeler.

\begin{cousinweb}
\textbf{\kn{ಚಕ್ರ}} (\emph{cakra}) is borrowed from Sanskrit \emph{cakra}, \textquotedblleft{}a wheel, a potter's wheel, a circle\textquotedblright{} (CDIAL 4538).
\end{cousinweb}

\begin{grammarlens}[title={What you've built}]
One more word for getting around.
\end{grammarlens}

\subsection*{Your turn}
\begin{itemize}
\raggedright
  \item \emph{Say it:} \emph{cakra}
  \item \emph{Say it:} \emph{cakra}, once more
  \item \emph{Say it:} \emph{saikalna hindina cakra ḍoṅkāgide.} — The back wheel of the bicycle is bent.
  \item \emph{Recall:} say the Kannada for a battery, then the Kannada for to ride, then say \emph{cakra} again
\end{itemize}

\begin{tcolorbox}[breakable,colback=teal!4,colframe=teal!35!black,title={Before you move on}]
What does \kn{ಚಕ್ರ} mean? (\textquotedblleft{}A wheel\textquotedblright{}.) Say it once more.
\end{tcolorbox}
\section[\texorpdfstring{\kn{ಟಯರ್} (ṭayar)}{ಟಯರ್ (ṭayar)}]{\kn{ಟಯರ್} (ṭayar) — a tyre}

\label{lesson:KA-C449-tayar}



\begin{quote}
Before the new one: say the Kannada for to ride, then the Kannada for a wheel.
\end{quote}

\subsection*{You'll want to know: \kn{ಟಯರ್}}
\textbf{\kn{ಟಯರ್}} — \emph{ṭayar} — \textquotedblleft{}a tyre\textquotedblright{}.

\begin{quote}
\kn{ಕಾರಿನ} \kn{ಟಯರ್} \kn{ಹಳೆಯದಾಗಿತ್ತು}.
\end{quote}

\begin{quote}
\emph{kārina ṭayar haḷeyadāgittu.}
\end{quote}

\begin{quote}
The car tyre was old.
\end{quote}

The rubber round the \kn{ಚಕ್ರ}.

\begin{cousinweb}
\textbf{\kn{ಟಯರ್}} (\emph{ṭayar}) is a loan from English \emph{tyre}.
\end{cousinweb}

\begin{grammarlens}[title={What you've built}]
One more word for getting around.
\end{grammarlens}

\subsection*{Your turn}
\begin{itemize}
\raggedright
  \item \emph{Say it:} \emph{ṭayar}
  \item \emph{Say it:} \emph{ṭayar}, once more
  \item \emph{Say it:} \emph{kārina ṭayar haḷeyadāgittu.} — The car tyre was old.
  \item \emph{Recall:} say the Kannada for to ride, then the Kannada for a wheel, then say \emph{ṭayar} again
\end{itemize}

\begin{tcolorbox}[breakable,colback=teal!4,colframe=teal!35!black,title={Before you move on}]
What does \kn{ಟಯರ್} mean? (\textquotedblleft{}A tyre\textquotedblright{}.) Say it once more.
\end{tcolorbox}
\section[\texorpdfstring{\kn{ಸ್ಟೆಪ್ನಿ} (sṭepni)}{ಸ್ಟೆಪ್ನಿ (sṭepni)}]{\kn{ಸ್ಟೆಪ್ನಿ} (sṭepni) — a spare tyre}

\label{lesson:KA-C449-stepni}



\begin{quote}
Before the new one: say the Kannada for a wheel, then the Kannada for a tyre.
\end{quote}

\subsection*{You'll want to know: \kn{ಸ್ಟೆಪ್ನಿ}}
\textbf{\kn{ಸ್ಟೆಪ್ನಿ}} — \emph{sṭepni} — \textquotedblleft{}a spare tyre\textquotedblright{}.

\begin{quote}
\kn{ನಮ್ಮ} \kn{ಕಾರಿನಲ್ಲಿ} \kn{ಸ್ಟೆಪ್ನಿ} \kn{ಇರಲಿಲ್ಲ}.
\end{quote}

\begin{quote}
\emph{namma kārinalli sṭepni iralilla.}
\end{quote}

\begin{quote}
There was no spare tyre in our car.
\end{quote}

Indian English for the spare wheel carried in a vehicle.

\begin{cousinweb}
\textbf{\kn{ಸ್ಟೆಪ್ನಿ}} (\emph{sṭepni}) is a loan from English \emph{stepney}, in Indian use, \textquotedblleft{}a spare wheel\textquotedblright{}.
\end{cousinweb}

\begin{grammarlens}[title={What you've built}]
One more word for getting around.
\end{grammarlens}

\subsection*{Your turn}
\begin{itemize}
\raggedright
  \item \emph{Say it:} \emph{sṭepni}
  \item \emph{Say it:} \emph{sṭepni}, once more
  \item \emph{Say it:} \emph{namma kārinalli sṭepni iralilla.} — There was no spare tyre in our car.
  \item \emph{Recall:} say the Kannada for a wheel, then the Kannada for a tyre, then say \emph{sṭepni} again
\end{itemize}

\begin{tcolorbox}[breakable,colback=teal!4,colframe=teal!35!black,title={Before you move on}]
What does \kn{ಸ್ಟೆಪ್ನಿ} mean? (\textquotedblleft{}A spare tyre\textquotedblright{}.) Say it once more.
\end{tcolorbox}
\section[\texorpdfstring{\kn{ಪಂಕ್ಚರ್} \kn{ಆಗು} (paṅkcar āgu)}{ಪಂಕ್ಚರ್ ಆಗು (paṅkcar āgu)}]{\kn{ಪಂಕ್ಚರ್} \kn{ಆಗು} (paṅkcar āgu) — to get a puncture}

\label{lesson:KA-C449-pankcar-agu}



\begin{quote}
Before the new one: say the Kannada for a tyre, then the Kannada for a spare tyre.
\end{quote}

\subsection*{You'll want to know: \kn{ಪಂಕ್ಚರ್} \kn{ಆಗು}}
\textbf{\kn{ಪಂಕ್ಚರ್} \kn{ಆಗು}} — \emph{paṅkcar āgu} — \textquotedblleft{}to get a puncture\textquotedblright{}.

\begin{quote}
\kn{ದಾರಿಯಲ್ಲಿ} \kn{ನನ್ನ} \kn{ಸ್ಕೂಟರ್ನ} \kn{ಟಯರ್} \kn{ಪಂಕ್ಚರ್} \kn{ಆಯಿತು}.
\end{quote}

\begin{quote}
\emph{dāriyalli nanna skūṭarna ṭayar paṅkcar āyitu.}
\end{quote}

\begin{quote}
On the way my scooter tyre got a puncture.
\end{quote}

Past (it …): \kn{ಪಂಕ್ಚರ್} \kn{ಆಯಿತು}, \emph{paṅkcar āyitu}.

The tyre is the subject: \kn{ಟಯರ್} \kn{ಪಂಕ್ಚರ್} \kn{ಆಯಿತು}.

\begin{cousinweb}
\textbf{\kn{ಪಂಕ್ಚರ್} \kn{ಆಗು}} (\emph{paṅcar āgu}) joins a loan from English \emph{puncture} and the native verb \kn{ಆಗು}, to become.
\end{cousinweb}

\begin{grammarlens}[title={What you've built}]
One more word for getting around.
\end{grammarlens}

\subsection*{Your turn}
\begin{itemize}
\raggedright
  \item \emph{Say it:} \emph{paṅkcar āgu}
  \item \emph{Say it:} \emph{paṅkcar āgu}, once more
  \item \emph{Say it:} \emph{dāriyalli nanna skūṭarna ṭayar paṅkcar āyitu.} — On the way my scooter tyre got a puncture.
  \item \emph{Recall:} say the Kannada for a tyre, then the Kannada for a spare tyre, then say \emph{paṅkcar āgu} again
\end{itemize}

\begin{tcolorbox}[breakable,colback=teal!4,colframe=teal!35!black,title={Before you move on}]
What does \kn{ಪಂಕ್ಚರ್} \kn{ಆಗು} mean? (\textquotedblleft{}To get a puncture\textquotedblright{}.) Say it once more.
\end{tcolorbox}
\section[\texorpdfstring{\kn{ಬ್ರೇಕ್} (brēk)}{ಬ್ರೇಕ್ (brēk)}]{\kn{ಬ್ರೇಕ್} (brēk) — a brake}

\label{lesson:KA-C449-brek}



\begin{quote}
Before the new one: say the Kannada for a spare tyre, then the Kannada for to get a puncture.
\end{quote}

\subsection*{You'll want to know: \kn{ಬ್ರೇಕ್}}
\textbf{\kn{ಬ್ರೇಕ್}} — \emph{brēk} — \textquotedblleft{}a brake\textquotedblright{}.

\begin{quote}
\kn{ಚಾಲಕ} \kn{ತಕ್ಷಣ} \kn{ಬ್ರೇಕ್} \kn{ಹಾಕಿದನು}.
\end{quote}

\begin{quote}
\emph{cālaka takṣaṇa brēk hākidanu.}
\end{quote}

\begin{quote}
The driver braked at once.
\end{quote}

\kn{ಬ್ರೇಕ್} \kn{ಹಾಕು}, to put on the brake.

\begin{cousinweb}
\textbf{\kn{ಬ್ರೇಕ್}} (\emph{brēk}) is a loan from English \emph{brake}.
\end{cousinweb}

\begin{grammarlens}[title={What you've built}]
One more word for getting around.
\end{grammarlens}

\subsection*{Your turn}
\begin{itemize}
\raggedright
  \item \emph{Say it:} \emph{brēk}
  \item \emph{Say it:} \emph{brēk}, once more
  \item \emph{Say it:} \emph{cālaka takṣaṇa brēk hākidanu.} — The driver braked at once.
  \item \emph{Recall:} say the Kannada for a spare tyre, then the Kannada for to get a puncture, then say \emph{brēk} again
\end{itemize}

\begin{tcolorbox}[breakable,colback=teal!4,colframe=teal!35!black,title={Before you move on}]
What does \kn{ಬ್ರೇಕ್} mean? (\textquotedblleft{}A brake\textquotedblright{}.) Say it once more.
\end{tcolorbox}
